\newpage
\begin{adjustwidth}{-5pt}{-2pt}
\begin{singlespace}
\tcbset{colback=white!10!white}
\begin{tcolorbox}[title=\textbf{\section{Model Card - \mfourttwo}\label{card:mfourttwo}},
    breakable, sharp corners, boxrule=0.5pt] 
\begin{mcsection}{Model Details\footnote{For this card, we use the template from \citet{10.1145/3287560.3287596}.}} 
\item Person or organization developing model: \textit{Developed by FAIR, Meta}
\item Model date: \textit{November 30, 2023}
\item Model version: \mfourtlgtwo 
\item Model type: \textit{Multitask-\unitytwo with (a) Conformer speech encoder, (b) Transformer text encoder-decoder and (c) Transformer encoder with a non-autoregressive decoder for \tu}.
\begin{itemize}
    \item \it The training algorithm of \mfourtlgtwo is described in the paper: \seamlesscitation 
\item License: \textit{CC-BY-NC 4.0
\footnote{\url{https://creativecommons.org/licenses/by-nc/4.0/legalcode}}}
\item Where to send questions or comments about the model: \\
\url{https://github.com/facebookresearch/seamless_communication/issues}
\end{itemize} 
\end{mcsection}


\begin{mcsection}{Intended Use}
    \item Primary intended uses: \textit{\mfourtlgtwo is a multilingual and multimodal translation model primarily intended for research in speech and text translation. It allows for:
    \begin{itemize}
        \item \asr: Automatic speech recognition for \nasrlangs languages.
        \item \sst: Speech-to-Speech translation from \nsslangs source speech languages into \ntslangs target speech languages. 
        \item \st: Speech-to-text translation from \nsslangs source speech languages into \ntextlangs target text languages.
        \item \ttst: Text-to-Speech translation from \ntextlangs source text languages into \ntslangs target speech languages.
        \item \mt: Text-to-text translation (MT) from \ntextlangs source text languages into \ntextlangs target text languages.
        \item TTS: Text-to-speech synthesis for 36 languages.
    \end{itemize}
Information on how to use the model can be found in the github repository at \url{https://github.com/facebookresearch/seamless_communication} along with scripts for evaluation and finetuning.}
\item Primary intended users: \textit{Primary users are researchers and machine translation (speech and text) research community. }
\item Out-of-scope use cases: \textit{\mfourtlgtwo is a research model and is not released for production deployment. \mfourtlgtwo is trained on general domain data and is not intended to be used with domain-specific inputs, such as the medical domain or legal domain. The model is not intended to be used for long-form translation. The model was trained on short text and speech inputs, therefore translating longer sequences might result in quality degradation. \mfourtlgtwo translations can not be used as certified translations.} 
\end{mcsection}

\begin{mcsection}{Metrics}
    \item Model performance measures: \textit{For the \st task, \mfourttwo models were evaluated using the \bleu metric adopted by SOTA models in speech-to-text translation. The models were additionally evaluated with \spbleu and \blaser on \st. For \sst, the models are evaluated with \asrbleu and \blaser. For the \mt taks, we report quality in terms of \chrf. For \asr, we report the widely adopted metric of \wer with the text normalized following the normalization in \cite{whisper}.
    Additionally, we performed human evaluations with the XSTS protocol, measured added toxicity, robustness, and bias, and reported red teaming results of \mfourtlgtwo.}
\end{mcsection}


\begin{mcsection}{Evaluation Data}
    \item Datasets: \textit{\fleurs, \flores, \covost and \cvss, \holisticbias and \multilingualholisticbias described in \citet{costa2023multilingual}.
    }
\item Motivation: \textit{We used \fleurs as it provides
an n-way parallel speech and text dataset in 102 languages, on which we can evaluate \mfourttwo models on multiple tasks.}

\end{mcsection}

\begin{mcsection}{Training Data}
\item \textit{We used parallel multilingual data from a variety of sources to train the model. For data statistics see \Cref{tbl:offline:s2tdata,tbl:offline:s2stdata} of \seamlesscitation
}
\end{mcsection}



\begin{mcsection}{Ethical Considerations}
 \item \textit{In this work, we took a comprehensive approach to prioritize human users and minimize risks that could be transferred to them. While we have documented various evaluation and responsible AI techniques deployed in our work, here are some additional points to highlight. For one, many languages chosen for this study are low-resource languages.
 While quality translation could improve world readiness and information access for many in these communities, such access could also make groups with lower levels of digital literacy more vulnerable to misinformation or online scams. The latter scenarios could arise if bad actors misappropriate our work for nefarious activities, which we conceive as an example of unintended use.
 Finally, although we did our best to optimize for translation quality, toxic, biased, or false outputs produced by the model could remain. These could have an adverse impact on those who rely on these translations to make important decisions (particularly when related to health and safety).} 
\end{mcsection}


\begin{mcsection}{Caveats and Recommendations}
  \item Limitations: \textit{Researchers should consider implementing additional integrity mitigations for ``added toxicity'' when using the model in a research application.}
\end{mcsection}


\end{tcolorbox}
\end{singlespace}
\end{adjustwidth}
